\section{What the System Does with More Intelligence}

An input does not acquire priority because it came from an intelligent machine. Once it enters a factory, a market, or an organization, the receiving system decides what happens next. Some inputs find a better route. Some accumulate. Some help drive the system unstable. These are different mechanisms, and they require different remedies.

\subsection{Rerouting: Moving Some of the Trial and Error}

The direct route is often the expensive route. A desired molecule is easy to describe and hard to synthesize. A production target is easy to enter into a spreadsheet and hard to reach with the equipment already installed. A search algorithm earns its keep when it discovers a path that asks less of the costly medium.

Call this a \emph{probabilistic shortcut}: spend computation to improve the odds of the next expensive intervention. Screen candidate compounds before synthesizing them. Simulate a warehouse layout before moving the racks. Search a design space before cutting metal. The saving comes from changing which physical trials are attempted, how they are ordered, and how much each trial teaches.

Monte Carlo methods made a related trade long before deep learning. Sampling replaces certain difficult calculations with estimates whose precision depends on the sampling scheme and budget. It does not make every estimate cheap or every rare event visible. It gives engineers a different way to buy an answer. In games, systems such as AlphaGo could search within rules that were already executable, with outcomes that could be obtained without building a physical board or waiting for a human opponent \cite{silver2016mastering}. The search environment supplied an unusually inexpensive judge.

The judge becomes more costly outside games. A predicted protein structure may narrow an experiment, suggest a mechanism, or support a screening workflow. The experimental question then matters: determining a structure, establishing binding, measuring activity, and demonstrating a therapeutic benefit are different tasks. A prediction need not be remeasured in every downstream use to be useful, but its uncertainty must be compatible with that use. Cheap structure proposals do not turn every biological claim into an established result.

Semiconductor fabrication illustrates another form of rerouting. When a direct scaling path becomes difficult, progress can continue through new materials, transistor structures, packaging, architecture, and specialized computation. Each route changes the cost distribution. More elaborate fabrication may buy performance at the price of more process steps, tighter tolerances, and expensive equipment. The result can be a real improvement without being a disappearance of resistance.

The same test applies to simulation. A simulator can fail a thousand times without breaking a gripper. Yet someone must establish which failures and successes transfer to the actual gripper. Domain randomization can improve transfer by exposing a policy to varied conditions; the range of variation remains a modeling decision \cite{tobin2017domain}. A shortcut pays when the physical trials it avoids cost more than the computation, modeling, and additional uncertainty it introduces.

The useful question is therefore concrete: \emph{which part of the real trial has been removed, which part has been moved, and which part still has to happen?} A project that cannot answer this may have relocated its bottleneck while mistaking the relocation for a breakthrough.

\subsection{Accumulation: Intelligence Becomes Inventory}

Consider an engineering team whose assistants can generate a hundred plausible patches a day. Its reviewers can investigate ten. Its test infrastructure can run more tests than before, but production incidents still require experienced staff. The organization now has a larger candidate inventory. Whether it has more reliable software depends on the receiving process.

For a queue with arrival rate $I(t)$ and absorption rate $A(S)$,
\begin{equation}
 \frac{\mathrm{d}S}{\mathrm{d}t}=I(t)-A(S),
 \qquad 0\leq A(S)\leq A_{\max}.
\end{equation}
$S$ counts unfinished items. The rates count items per unit time. At an empty queue, departures cannot exceed available work; a more detailed model would also include rejection, abandonment, rework, and different classes of item. The accounting point survives: sustained arrivals above effective service capacity create a growing backlog.

A backlog is stored time. It can be useful. Hospitals hold spare supplies, cloud systems buffer bursts, and laboratories retain promising candidates until an instrument becomes available. Inventory becomes a problem when it conceals an absorption limit. A long list of unreviewed proposals may look like a rich pipeline while making urgent work harder to identify.

Even stable queues can have long delays when utilization is high and arrivals or service times vary. Spare capacity buys something that an average throughput chart can hide: the ability to absorb a surprise. Running every expert, machine, and approval channel at nominal capacity leaves little room to recover from an exception.

Accumulation also moves upstream. If an experimental facility is booked months ahead, researchers change what they propose. If deployment review becomes the bottleneck, developers optimize for the review process. The downstream constraint starts shaping the supposedly unconstrained search space.

Supply chains show how this can become oscillation. Demand signals pass through ordering policies, inventory targets, production schedules, and delivery delays. Local adjustments can amplify a small change in final demand into much larger upstream orders: the bullwhip effect \cite{lee1997information}. Better information can help, but information alone does not remove lead times or correct a badly tuned replenishment rule.

There is no universal clock assigning short cycles to light assets and long cycles to heavy ones. Nevertheless, adjustment times matter. Clearing a content queue and rebuilding a power plant operate on different clocks. A decision that was appropriate when initiated may reach completion in a changed environment. Heavy commitments preserve old expectations in physical form.

Reducing arrivals can relieve a queue; increasing service capacity can raise its sustainable throughput. Improving candidate quality can do both by reducing rework and increasing the useful output of each review. These are distinct interventions. Generating more candidates addresses none of them by itself.

\subsection{Instability: Correction Becomes a Disturbance}

Oscillation is not the inevitable outcome of fast computation. A well-designed controller can use rapid measurements to stabilize a slow plant. The danger is a control policy whose gain, delay, and model do not fit the plant it is controlling.

An operator sees a shortfall and orders more capacity. Before the capacity arrives, the shortfall persists, so the operator orders still more. When the first expansion comes online, the later orders remain in motion. The controller has interpreted the absence of an immediate response as evidence that its previous action was too weak.

This is how a sensible local instruction can create a damaging global rhythm. A warehouse keeps rerouting staff before any team has settled into its new task. A company changes quarterly objectives before the previous reorganization has produced interpretable evidence. An automated market strategy reacts to liquidity that disappears as other strategies react to it.

The official investigation of the May 2010 US market disruption documented interactions among a large sell program, changing liquidity, and trading behavior \cite{sec2010findings}. It is a historical case of interacting mechanisms, rather than proof that trading speed alone causes crashes. A fast system needs an account of what its own action does to the signals it will read next.

Protective devices enter that same loop. A trading halt can create time for liquidity and information to recover. A known threshold can also alter behavior before the halt: participants may trade earlier because they expect access to close. Theoretical work on price limits describes such a magnet effect \cite{subrahmanyam1994}. A threshold must be evaluated with the behavior it induces.

In physical control, inertia, damping, and active compensation have distinct roles. Inertia slows a change in velocity or frequency; damping dissipates an oscillation; a controller can supply a response that depends on its design and available energy. Calling all three “friction” obscures the engineering. Their common relevance here is that rapid decisions need a receiving system capable of turning them into a stable trajectory.

Three questions separate the interventions. Can search identify a route that needs fewer costly trials? Can the receiving process absorb the proposed work? Does the feedback policy remain stable with actual delays and disturbances? The corresponding investments are better validation, better absorption, and better control.

More intelligence can improve all three. It can also flood a review queue, accelerate a badly chosen intervention, or exploit a shortcut that works only inside the model. The difference lies in what becomes of the input after it crosses the interface.
